\section{Repository audit and the common denominator}
\label{sec:audit}

\subsection{The checked base and the research layer}

The repository has a useful but important split.  The \code{ConciseIPRA}
subtree implements a scoped first-order syntax, intuitionistic natural
deduction with equality, the ordinary presentation of iPRA, and translations
to a compressed native scalar calculus.  In particular, a theorem of type
\[
  \Gamma\vdash[\PRAAx]\varphi
\]
is an actual derivation in the object theory, while
\code{pra\_iff\_native} transports such a derivation to the native calculus.
The direction of this bridge matters: it does not turn an arbitrary semantic
Lean theorem into an iPRA proof.

The \code{docs/proofs} subtree is a broad research laboratory.  It includes
packets on exact algebra, roots and branches, constructive analysis, geometric
certificates, finite probability, arithmetic geometry, analytic number theory,
and several proposed foundations.  Many packets have polished mathematical
reports and executable finite prototypes.  Some include checked semantic Lean
lemmas; a smaller number contain combinators whose conclusions really have
\code{Prf PRAAx} type.  Most large application claims, however, stop before a
complete object-language derivation.

This report treats that heterogeneity as data rather than as a defect.  The
packets expose a repeated architecture, and the missing object proof is usually
located at one of a small number of interfaces.

\subsection{Recurring successful patterns}

Across the packets, the following patterns repeatedly work.

\begin{enumerate}[label=\textbf{P\arabic*.}]
\item \textbf{Finite coding.}  Objects are represented by natural-number codes:
      bounded lists, dyadic rationals, finite matrices, polynomials, boxes,
      graphs, or certificate traces.
\item \textbf{Boolean verification.}  A fixed primitive-recursive checker
      verifies all finite side conditions.  Bounded quantifiers are folded into
      finite conjunctions, disjunctions, and searches.
\item \textbf{Quantifier-free invariant.}  A bounded computation is justified
      by an equation-valued invariant; quantified assumptions stay in the
      context and are instantiated only at the finitely many queried terms.
\item \textbf{Guarded arithmetic.}  Division, normalization, root labels,
      chart choices, and inverse functions are used only with explicit rational
      margins.
\item \textbf{Relational output.}  Near singularities or collisions, the output
      is a finite relation, cluster, multiset, interval, or cover rather than a
      discontinuously selected exact witness.
\item \textbf{Certificate separation.}  A small verifier is distinguished from
      the possibly deep theorem or algorithm that produces an accepted
      certificate.
\end{enumerate}

These patterns are already sufficient for a large constructive fragment of
ordinary mathematics.  Their weakness is not mathematical but architectural:
several packets reimplement the same finite mechanisms, and the interface from
the mechanism to \code{Prf PRAAx} is often left implicit.

\subsection{Recurring gaps}

The audit groups the reported difficulties into six classes.

\paragraph{G1: semantic/object confusion.}
A theorem about a host-language function, or a passing numerical test, is used
as if it were a derivation in iPRA.  The remedy is to make every local premise
and every exported endpoint explicitly typed as \(\Prf\,\PRAAx\).

\paragraph{G2: wrong induction shape.}
The intended invariant contains ``for every precision'', an unbounded
existential, or a completed real-number assertion.  iPRA permits induction only
for quantifier-free formulas.  The repair is to fix the requested precision,
compile all bounded logic to Boolean equations, run the finite proof, and only
then introduce universal quantifiers.

\paragraph{G3: hidden infinite choice.}
For each precision there exists a finite witness, but separately chosen
witnesses may be incompatible.  Backward viability solves fixed finite
horizons.  Coherent infinite names require either a direct pairwise stability
estimate or a shrinking viable-cloud theorem with a modulus.

\paragraph{G4: discontinuous selection.}
Ordered roots at collisions, signs at zero, normalized axes at the identity,
and branches at singular points cannot be chosen uniformly from approximate
input.  The stable object must be changed, or a positive gap must be supplied.

\paragraph{G5: ineffective producer.}
Finiteness, existence, or a cardinality bound does not provide a search cutoff,
height bound, zero-free constant, rational covering representative, or
convergence rate.  A verifier can still be completely formalized, but the
producer remains a separate theorem.

\paragraph{G6: code extensionality.}
Different codes may represent the same object.  Operations must respect an
observational equivalence or a finite back-and-forth relation; raw code equality
is generally too strong.

\subsection{A single target architecture}

The common denominator is a stack of four layers.

\begin{center}
\begin{tabularx}{0.94\textwidth}{@{}lX@{}}
\toprule
Layer & Responsibility \\
\midrule
Representation & Fixed natural-number codes and fixed \(\PR\) symbols. \\
Finite proof kernel & Boolean folds, bounded search, traces, demands, paths,
                      viability, rational error arithmetic. \\
Mathematical adapter & The theorem-specific local inequalities, algebraic
                       identities, separation estimates, or certificate
                       soundness lemmas. \\
Object export & A value of type \(\Prf\,\PRAAx\), followed only then by the
                checked native translation. \\
\bottomrule
\end{tabularx}
\end{center}

The framework developed below makes the second layer reusable and forces the
third layer to state its genuine missing obligations.  It does not pretend to
solve an ineffective producer problem by changing notation.
